% Resumo em língua portuguesa
\setlength{\absparsep}{18pt}
\begin{resumo}
Os Eventos Adversos (EAs) em ambiente hospitalar configuram um dos maiores desafios contemporâneos da saúde pública global, figurando entre as principais causas evitáveis de morbimortalidade em pacientes internados. Historicamente, os sistemas de vigilância hospitalar baseiam-se em notificações voluntárias e espontâneas, modelo sabidamente ineficaz que captura menos de 10\% dos incidentes com dano real, perpetuando uma falsa sensação de segurança institucional. Em contrapartida, a metodologia \textit{Global Trigger Tool} (GTT), desenvolvida pelo \textit{Institute for Healthcare Improvement} (IHI), consolidou-se internacionalmente como o padrão-ouro para rastreamento retrospectivo de EAs mediante a detecção de gatilhos clínicos e classificação padronizada de gravidade segundo o \textit{National Coordinating Council for Medication Error Reporting and Prevention} (NCC MERP). No entanto, o ensino formal da segurança do paciente em cursos de graduação em saúde enfrenta graves entraves bioéticos, riscos jurídicos e limitações de privacidade impostas pela Lei Geral de Proteção de Dados (LGPD), inviabilizando o acesso irrestrito de discentes a prontuários médicos reais. Este projeto departamental apresenta o desenvolvimento e a formalização técnica do \textbf{SIGEA-GTT} (Sistema Inteligente de Gestão de Eventos Adversos – Global Trigger Tool), um ambiente web inovador voltado ao treinamento, simulação e auditoria clínica no âmbito da Universidade Federal de Sergipe (UFS). O sistema opera com isolamento estrito de dados sensíveis, provendo casos clínicos fictícios estruturados (notas multidisciplinares, prescrições aprazadas, exames laboratoriais e procedimentos cirúrgicos), cronômetro de auditoria de 20 minutos e rastreamento dos 53 gatilhos canônicos do IHI divididos em 6 módulos assistenciais. Como diferencial pedagógico, o SIGEA-GTT integra nativamente seis Ferramentas da Qualidade (Diagrama de Ishikawa 6M, Matriz GUT com cálculo em tempo real, Matriz 5W2H, Ciclo PDCA/PDSA, Matriz SWOT e Brainstorming), além de avaliação formativa com parecer do docente e cálculo automático dos três indicadores epidemiológicos canônicos do IHI: taxa de EAs por 1.000 pacientes-dia, taxa de EAs por 100 admissões e percentual de pacientes com EAs. A solução foi arquitetada em Java 21 LTS com Spring Boot 3.3, Angular 18+ com componentes \textit{Standalone} e reatividade baseada em \textit{Signals}, persistência relacional em PostgreSQL 16 com campos JSONB e segurança robusta com autenticação JWT \textit{stateless} restrita ao domínio institucional acadêmico. Os resultados alcançados consolidam uma plataforma pioneira no Brasil para o ensino rigoroso de segurança do paciente fundamentado em métodos de engenharia de software de alta precisão.

\textbf{Palavras-chave}: Segurança do Paciente. Global Trigger Tool. Eventos Adversos. Ferramentas da Qualidade. Simulação Clínica Computacional. Engenharia de Software.
\end{resumo}